% ============================================================================
% KATAPENGANTAR.TEX
% Kata Pengantar Buku Ajar
% ============================================================================

\chapter*{Kata Pengantar}
\addcontentsline{toc}{chapter}{Kata Pengantar}

\vspace{0.5cm}

Puji syukur kehadirat Allah SWT yang telah memberikan rahmat dan karunia-Nya sehingga penulis dapat menyelesaikan buku ajar \textbf{Pemrograman Game} ini. Buku ajar ini disusun sebagai bahan pembelajaran untuk mata kuliah Pemrograman Game di Program Studi S1 Informatika dengan menggunakan pendekatan Outcome-Based Education (OBE).

Buku ajar ini membahas pengembangan game menggunakan Unity Engine dan bahasa pemrograman C\#. Materi disusun secara sistematis mulai dari konsep dasar game development, pengenalan Unity Engine, pemrograman C\# untuk Unity, hingga aspek-aspek teknis pengembangan game seperti fisika, animasi, UI/UX, audio, AI, optimasi, dan deployment.

\textbf{Keunikan buku ajar ini} dibandingkan dengan buku-buku sejenis adalah:
\begin{itemize}
\item Menggunakan pendekatan Outcome-Based Education (OBE) dengan pemetaan CPL-CPMK-Bab-Section yang jelas
\item Fokus pada pembelajaran berbasis proyek (project-based learning) yang dominan praktik
\item Menggunakan sumber referensi open-access yang dapat diakses secara bebas
\item Disusun sesuai dengan Rencana Pembelajaran Semester (RPS) 16 pertemuan
\item Setiap bab dilengkapi dengan learning outcomes, peta konsep, ringkasan teori inti, latihan, dan ringkasan
\end{itemize}

\textbf{Sasaran pembaca} buku ajar ini adalah:
\begin{itemize}
\item Mahasiswa Program Studi S1 Informatika yang mengambil mata kuliah Pemrograman Game
\item Dosen pengampu mata kuliah Pemrograman Game sebagai bahan referensi pengajaran
\item Praktisi dan pengembang game pemula yang ingin mempelajari Unity Engine
\item Pembaca umum yang tertarik dengan pengembangan game menggunakan Unity
\end{itemize}

Penulis menyadari bahwa buku ajar ini masih jauh dari sempurna. Oleh karena itu, kritik dan saran yang membangun sangat diharapkan untuk perbaikan edisi selanjutnya.

Akhir kata, penulis mengucapkan terima kasih kepada:
\begin{itemize}
\item Program Studi S1 Informatika yang telah memberikan kesempatan untuk menyusun buku ajar ini
\item Semua pihak yang telah memberikan dukungan dan masukan dalam penyusunan buku ajar ini
\item Unity Technologies dan komunitas Unity yang telah menyediakan dokumentasi dan resources yang sangat membantu
\item Semua sumber referensi open-access yang telah digunakan sebagai acuan dalam penulisan buku ajar ini
\end{itemize}

Semoga buku ajar ini dapat memberikan manfaat bagi pembaca dalam mempelajari pengembangan game menggunakan Unity Engine.

\vfill

\begin{flushright}
\textbf{[Nama Kota], [Bulan] [Tahun]}\\[0.5cm]
\textbf{Penulis}
\end{flushright}
